\documentclass[11pt]{article}
\usepackage[margin=1in]{geometry}
\usepackage{enumitem}
% =============================================================================
% Preambles/header.tex — shared LaTeX/Beamer preamble for this project
%
% Use from a Beamer lecture by prepending:
%     \input{header}
% and compiling with TEXINPUTS=../Preambles:$TEXINPUTS (the /compile-latex
% skill does this automatically).
%
% PALETTE CONTRACT. Color names defined here MUST match the names in
% Quarto/theme-template.scss so Beamer and Quarto renderings use the same
% palette. The scripts/check-palette-sync.sh script enforces this.
%
% To customize: edit the HEX values below AND the matching $variable in
% Quarto/theme-template.scss, then run scripts/check-palette-sync.sh.
% =============================================================================

% -----------------------------------------------------------------------------
% Engine detection + encoding
%
% Our /compile-latex skill uses XeLaTeX, where inputenc is unnecessary and
% can produce encoding warnings. Only load inputenc under pdfTeX.
% -----------------------------------------------------------------------------
\usepackage{iftex}
\ifPDFTeX
  \usepackage[utf8]{inputenc}
\fi

\usepackage{amsmath, amssymb, amsthm}
\usepackage{booktabs}
\usepackage{graphicx}

% -----------------------------------------------------------------------------
% PALETTE — mirrors Quarto/theme-template.scss variable names.
% Load xcolor and define colors BEFORE anything that references them
% (notably hyperref, hypersetup, and the Beamer color assignments).
% -----------------------------------------------------------------------------
\usepackage{xcolor}

\definecolor{primary-blue}{HTML}{012169}      % headings, accents
\definecolor{primary-gold}{HTML}{B9975B}      % emphasis, borders
\definecolor{highlight-yellow}{HTML}{F2A900}  % markers, alerts
\definecolor{light-bg}{HTML}{E8EDF5}          % title / section backgrounds
\definecolor{jet}{HTML}{1A1A1A}               % body text

% Semantic colors (used in diagrams and annotations)
\definecolor{positive}{HTML}{15803D}          % good / observed / identified
\definecolor{negative}{HTML}{B91C1C}          % bad / problematic / bias
\definecolor{neutral}{HTML}{525252}           % reference / context

% Highlight accents (match .hi-* classes in the SCSS)
\definecolor{hi-slate}{HTML}{314F4F}
\definecolor{hi-green}{HTML}{2E8B57}
\definecolor{hi-red}{HTML}{C41E3A}

% Semantic aliases so skill templates and snippets can write
% \textcolor{accent}{...} without hard-coding "primary-blue".
\colorlet{accent}{primary-blue}
\colorlet{accent2}{primary-gold}

% -----------------------------------------------------------------------------
% Hyperref — loaded AFTER xcolor + palette so link colors resolve.
% -----------------------------------------------------------------------------
\usepackage{hyperref}
\hypersetup{colorlinks=true, linkcolor=primary-blue, urlcolor=primary-blue,
            citecolor=primary-blue, pdfborder={0 0 0}}

% -----------------------------------------------------------------------------
% Course code — set once per deck (after \input{header}, before \begin{document})
% via \coursecode{CS401}. Shown in the footer on every slide so a deck's course
% is identifiable without opening the title page. Empty by default.
% -----------------------------------------------------------------------------
\makeatletter
\newcommand{\coursecode}[1]{\gdef\@coursecodetext{#1}}
\gdef\@coursecodetext{}
\makeatother

% -----------------------------------------------------------------------------
% Beamer theme assignments (only apply under Beamer).
%
% We detect Beamer via \@ifclassloaded{beamer}, which is the documented,
% non-internal way. This must be wrapped in \makeatletter/\makeatother
% because \@ifclassloaded contains an @-catcode character.
% -----------------------------------------------------------------------------
\makeatletter
\@ifclassloaded{beamer}{%
  \usefonttheme{professionalfonts}
  \setbeamercolor{structure}{fg=primary-blue}
  \setbeamercolor{frametitle}{fg=primary-blue}
  \setbeamercolor{title}{fg=primary-blue}
  \setbeamercolor{subtitle}{fg=primary-gold}
  \setbeamercolor{author}{fg=primary-blue}
  \setbeamercolor{institute}{fg=neutral}
  \setbeamercolor{date}{fg=primary-gold}
  \setbeamercolor{itemize item}{fg=primary-blue}
  \setbeamercolor{itemize subitem}{fg=primary-gold}
  \setbeamercolor{alerted text}{fg=primary-gold}
  \setbeamercolor{block title}{fg=white, bg=primary-blue}
  \setbeamercolor{block body}{bg=light-bg}
  % Semantic block variants: block=definition/concept (blue), exampleblock=
  % worked example (green/positive), alertblock=key takeaway or caution (gold).
  % Keep to <=2 colored boxes per slide across ALL three variants (INV-7).
  \setbeamercolor{block title example}{fg=white, bg=positive}
  \setbeamercolor{block body example}{bg=light-bg}
  \setbeamercolor{block title alerted}{fg=white, bg=primary-gold}
  \setbeamercolor{block body alerted}{bg=light-bg}

  % Minimal footer: course code (left) + page X / N (right), no navigation clutter
  \setbeamertemplate{navigation symbols}{}
  \setbeamertemplate{footline}{%
    \color{neutral}\footnotesize\hspace{0.7em}\@coursecodetext\hfill
    \insertframenumber/\inserttotalframenumber\hspace{0.7em}\vspace{0.3em}}
}{}
\makeatother

% -----------------------------------------------------------------------------
% TikZ setup — libraries the snippet gallery uses
% -----------------------------------------------------------------------------
\usepackage{tikz}
\usetikzlibrary{arrows.meta, positioning, calc, decorations.pathreplacing,
                fit, shapes.geometric, backgrounds}

% Shared TikZ styles — snippets and hand-written diagrams can pick these up
% via `\begin{tikzpicture}[every node/.style={...}]` or by naming specific
% styles (e.g., `\node[dag-node] (X) at (0,0) {X};`).
\tikzset{
  % Node styles for causal DAGs and similar
  dag-node/.style = {
    draw=primary-blue, fill=light-bg, rounded corners=2pt,
    minimum width=1.2cm, minimum height=0.9cm, align=center, font=\small
  },
  decision-node/.style = {
    draw=primary-gold, fill=white, diamond, aspect=1.8,
    minimum width=1.8cm, minimum height=1.2cm, align=center, font=\small,
    inner sep=1pt
  },
  % Edge styles — observed vs counterfactual
  observed-edge/.style = {-{Latex[length=2.5mm]}, thick, draw=jet},
  counterfactual-edge/.style = {-{Latex[length=2.5mm]}, thick, draw=neutral, dashed},
  confound-edge/.style = {-{Latex[length=2.5mm]}, thick, draw=negative, dashed},
  % Data-point styles for plots
  observed-dot/.style = {fill=primary-blue, draw=primary-blue, circle, inner sep=1.2pt},
  counterfactual-dot/.style = {fill=white, draw=neutral, circle, inner sep=1.2pt}
}

% -----------------------------------------------------------------------------
% Convenience macros
% -----------------------------------------------------------------------------
\newcommand{\muted}[1]{\textcolor{neutral}{#1}}
\newcommand{\key}[1]{\textcolor{primary-gold}{\textbf{#1}}}
\newcommand{\good}[1]{\textcolor{positive}{#1}}
\newcommand{\bad}[1]{\textcolor{negative}{#1}}

% Standout frame macro: full-bleed dark background for section transitions.
% Beamer-only (uses \begin{frame}/\setbeamercolor/\usebeamerfont) -- do not
% call from an article-class file (e.g. Notes/<CODE>/*-notes.tex). \muted,
% \key, \good, \bad, and \coursecode above are plain xcolor/gdef and are
% safe to use from either class; verified when Notes/article-class support
% was added.
% Expand: \transitionslide{Your transition title}
\newcommand{\transitionslide}[1]{%
  \begingroup
  \setbeamercolor{background canvas}{bg=primary-blue}
  \begin{frame}[plain]
    \centering
    \vfill
    {\usebeamerfont{title}\color{white}\Large #1\par}
    \vfill
  \end{frame}
  \endgroup
}

% Section divider frame (Beamer-only), an alternative to \transitionslide with a
% darker two-line style: near-black (jet) background, a small white label line
% above a larger gold title, and a thin gold rule along the bottom edge.
% Expand: \sectiondivider{Section 2}{Pointers}
\newcommand{\sectiondivider}[2]{%
  \begingroup
  \setbeamercolor{background canvas}{bg=jet}
  \begin{frame}[plain]
    \vspace*{3.0cm}
    {\color{white}\ttfamily\large #1\par}
    \vspace{0.5cm}
    {\color{primary-gold}\LARGE #2\par}
    \begin{tikzpicture}[remember picture, overlay]
      \draw[primary-gold, line width=2.5pt]
        ($(current page.south west)+(0,0.1)$) -- ($(current page.south east)+(0,0.1)$);
    \end{tikzpicture}
  \end{frame}
  \endgroup
}

% -----------------------------------------------------------------------------
% Channel watermark — "Concepts That Click" (YouTube).
%
% Article-class files (CompetitiveExam Q/A sets, Notes, Assignments) get a
% light diagonal draft watermark on every page plus a centered footer naming
% the channel. Beamer decks get a subtle TikZ background watermark instead
% (the footline already carries the course code). Centralized here so every
% MCQ PDF is branded without per-file edits.
% -----------------------------------------------------------------------------
\makeatletter
\@ifclassloaded{article}{%
  \usepackage{draftwatermark}
  \SetWatermarkText{Concepts That Click}
  \SetWatermarkScale{1}
  \SetWatermarkFontSize{2cm}
  \SetWatermarkLightness{0.86}
  \SetWatermarkAngle{45}
  \usepackage{fancyhdr}
  \pagestyle{fancy}
  \fancyhf{}
  \fancyfoot[C]{\footnotesize\textcolor{neutral}{Concepts That Click~|~YouTube: \href{https://www.youtube.com/@concepts.thatclick}{@concepts.thatclick}}}
  \renewcommand{\headrulewidth}{0pt}
  \renewcommand{\footrulewidth}{0pt}
  \fancypagestyle{plain}{%
    \fancyhf{}%
    \fancyfoot[C]{\footnotesize\textcolor{neutral}{Concepts That Click~|~YouTube: \href{https://www.youtube.com/@concepts.thatclick}{@concepts.thatclick}}}%
    \renewcommand{\headrulewidth}{0pt}%
    \renewcommand{\footrulewidth}{0pt}%
  }
}{}
\@ifclassloaded{beamer}{%
  \setbeamertemplate{background}{%
    \begin{tikzpicture}[remember picture, overlay]
      \node[rotate=45, opacity=0.055, align=center]
        at (current page.center)
        {\Huge\textcolor{neutral}{Concepts That Click}};
    \end{tikzpicture}%
  }
}{}
\makeatother 
\coursecode{JKSSB}
\renewcommand{\thesection}{07.\arabic{section}}
\newtheorem{example}{Example}[section]
\newenvironment{definitionbox}[1]{%
  \par\noindent\textbf{Definition (#1).}\ \itshape
}{\par\vspace{0.3em}}
\title{Computer Memory Hierarchy --- Detailed Lecture Notes}
\author{JKSSB Computer Awareness}
\date{\today}

\begin{document}
\maketitle

\section{Why a memory hierarchy exists}
A running computer never keeps all its data and instructions in one place.
Instead it uses several stores arranged in a hierarchy. The reason is a
three-way trade-off: the memory that is fastest is also the smallest and the
most expensive per bit, while the memory that is largest and cheapest is the
slowest. No single memory can be fast, large, and cheap at the same time.
The hierarchy places the small, fast, costly memories close to the CPU and
the large, slow, cheap memories further away.

\begin{definitionbox}{Memory hierarchy}
The ordered arrangement of storage levels in a computer, from the fastest
and smallest (registers) through cache and main memory down to the slowest
and largest (secondary and offline storage), chosen to balance speed, cost,
and capacity.
\end{definitionbox}

\begin{center}
\begin{tabular}{p{0.24\textwidth}p{0.66\textwidth}}
\toprule
\textbf{Term} & \textbf{Meaning in this lesson} \\
\midrule
Register & The fastest and smallest storage, located inside the CPU. \\
Cache & A small, very fast memory that holds recently and frequently used data. \\
Main memory (RAM) & The primary memory that holds active programs and data; it is volatile. \\
ROM & Non-volatile read-only primary memory that holds firmware such as the BIOS. \\
Secondary storage & Non-volatile storage such as an SSD or HDD, used for files. \\
Tertiary / offline & Backup and archival storage such as magnetic tape or cloud backup. \\
\bottomrule
\end{tabular}
\end{center}

The rest of this lesson uses these terms in one consistent sense. A
\textbf{register} is the CPU's own storage. \textbf{Cache} is the small fast
memory between the CPU and main memory. \textbf{Main memory} means RAM: the
primary memory that holds what is running. \textbf{Secondary storage} means
SSD and HDD, which keep files without power. \textbf{Tertiary or offline}
storage is the slowest level, used for backup and archive.

\section{The ordered hierarchy}
The levels, read from fastest to slowest, are registers, cache, main memory
(RAM), secondary storage, and then tertiary or offline storage.

\begin{center}
\begin{tabular}{p{0.20\textwidth}p{0.18\textwidth}p{0.16\textwidth}p{0.36\textwidth}}
\toprule
\textbf{Level} & \textbf{Speed} & \textbf{Size} & \textbf{Typical role} \\
\midrule
Registers & Fastest & Smallest & Values in use right now \\
Cache (L1/L2/L3) & Very fast & Small & Recently and frequently used data \\
Main memory (RAM) & Moderately fast & Larger & Active programs and data \\
Secondary (SSD/HDD) & Slower & Large & Permanent file storage \\
Tertiary/offline (tape, cloud) & Slowest & Largest & Backup and archive \\
\bottomrule
\end{tabular}
\end{center}

Reading downward, capacity tends to grow and speed tends to fall. The
anchor to remember is at the top: \textbf{registers are the fastest} storage
in the computer. Below them cache is still very fast but slower than the
registers. Main memory (RAM) is larger than cache but slower. Secondary
storage (SSD or HDD) is larger again, cheaper per bit, and slower, but it
keeps its contents without power. Tertiary or offline storage (magnetic
tape, cloud backup) is the slowest and cheapest level and is used for backup
and long-term archive.

\section{Registers and cache}
Registers are located inside the CPU and are the fastest and smallest
storage in the machine. Because they are so small, they hold only the values
needed at this instant: instructions, addresses, data, and intermediate
results.

Cache is a small, very fast memory that sits between the CPU and main
memory. It holds recently and frequently used data and instructions, and it
is faster than main memory. Cache does not store new information of its own;
it keeps \emph{copies} of data that really live in RAM. It is therefore not
a backup, and it loses its contents when power is removed, just as RAM does.

\begin{definitionbox}{Cache}
A small, very fast memory that holds recently and frequently used data and
instructions so the CPU can avoid the slower trip to main memory.
\end{definitionbox}

Cache is usually described in levels. \textbf{L1} is the smallest and
fastest and sits closest to the CPU core. \textbf{L2} is larger and a little
slower than L1. \textbf{L3} is the largest and slowest of the three, is
often shared across cores, and is still faster than main memory. For exam
purposes, the key point is the ordering: L1 is fastest and smallest, and L3
is largest and slowest among the three, but all cache is faster than RAM.

\begin{example}
While a student edits a document on the office desktop, the same lines and
the nearby rows are read again and again. The cache keeps those recently
used and nearby values ready, so the CPU does not have to wait for main
memory on every access. The cache holds copies of those values; the master
copy remains in RAM.
\end{example}

\section{Main memory (RAM) and ROM}
Main memory is the computer's primary memory. The programs and data that
are currently in use are held here while the machine runs, and the CPU works
directly with main memory. \textbf{RAM} (Random Access Memory) is
\emph{volatile}: its contents are lost when the power is switched off. RAM
is larger than cache but slower.

\textbf{ROM} (Read-Only Memory) is the other part of primary memory. It is
\emph{non-volatile}: it keeps its contents without power. In normal
operation it can be read but not written; changing it needs a special
process. ROM holds firmware such as the startup (BIOS) instructions.

\begin{definitionbox}{Volatile memory}
Memory that loses its contents when the power is removed. RAM is the common
example.
\end{definitionbox}

\begin{definitionbox}{Non-volatile memory}
Memory that retains its contents without power. ROM, SSD, and HDD are common
examples.
\end{definitionbox}

\begin{example}
An item states, ``RAM is non-volatile and will keep the file after
shutdown.'' This is wrong. RAM is volatile; its contents are lost when the
power goes off. It is \textbf{ROM} that is non-volatile. The wrong option has
swapped the two primary-memory types.
\end{example}

\section{Primary memory vs secondary storage}
\textbf{Primary memory} is directly accessible by the CPU and holds the
active programs and data. It includes RAM, which is volatile, and ROM, which
is non-volatile. \textbf{Secondary storage}, such as an SSD or HDD, is
non-volatile and keeps files and programs between sessions. It is larger and
cheaper per bit but slower than primary memory. The CPU does not access
secondary storage directly; data moves through main memory under the control
of the operating system.

\section{Secondary and tertiary storage}
Secondary storage comes in two common forms. A \textbf{solid-state drive
(SSD)} uses flash memory and has no moving parts; a \textbf{hard disk drive
(HDD)} uses rotating magnetic disks. Both are non-volatile and keep data
without power. \textbf{Tertiary or offline} storage, such as magnetic tape
and cloud backup, is the slowest and cheapest level and is used for backup
and long-term archive.

\section{Locality of reference}
Programs do not access memory at random; they tend to reuse the same data
and instructions. This behaviour is called \textbf{locality of reference},
and it comes in two forms. \textbf{Temporal locality} means the same item is
accessed again soon after it was last used, such as a loop counter used on
every iteration. \textbf{Spatial locality} means items stored at nearby
addresses are accessed next, such as the consecutive elements of an array.
This reuse is exactly why cache works: by keeping recently used and nearby
data ready, cache reduces the average time the CPU waits for data.

\begin{definitionbox}{Locality of reference}
The tendency of a running program to reuse recently accessed data
(temporal locality) and data at nearby addresses (spatial locality).
\end{definitionbox}

\section{Exam retrieval and common confusions}
Five distinctions prevent most errors on this topic.
\begin{enumerate}
  \item RAM is \textbf{volatile}; ROM is \textbf{non-volatile}. Do not swap
    them (TRAP-01).
  \item Cache is \textbf{not} a backup. It holds copies of data that live in
    RAM.
  \item RAM is main memory, \textbf{not} storage for files. Files belong in
    secondary storage.
  \item Cache is \textbf{faster} than main memory; reversing this is a common
    wrong answer (PYQ-03).
  \item The order of speed is registers, then cache, then RAM, then
    secondary storage. Registers are fastest.
\end{enumerate}

\begin{example}
An item asks which memory is fastest. Trace the hierarchy from the top:
registers are fastest and smallest, then cache, then main memory (RAM), then
secondary storage. The answer is registers. If the stem instead asks which
level holds recently used data, the answer is cache, and if it asks which is
volatile primary memory, the answer is RAM.
\end{example}

\subsection*{Exercises}
\begin{enumerate}[label=\arabic*.]
  \item List the levels of the memory hierarchy from fastest to slowest.
  \item Explain why the faster memories of the hierarchy are placed close to
    the CPU.
  \item Distinguish volatile memory from non-volatile memory, giving one
    example of each.
  \item Why is cache described as faster than main memory, and why is it not
    a backup?
  \item Define temporal locality and spatial locality, and explain how they
    relate to cache.
\end{enumerate}

\subsection*{Solutions}
\begingroup\small
\begin{enumerate}[label=\arabic*.]
  \item Registers, cache (L1, L2, L3), main memory (RAM), secondary storage
    (SSD/HDD), then tertiary/offline storage.
  \item Faster memory is smaller and more expensive per bit. Placing it close
    to the CPU keeps the values the CPU needs soonest at the fastest level,
    so the average access time is low while cost and capacity are balanced.
  \item Volatile memory loses its contents when power is removed --- RAM is
    the example. Non-volatile memory retains its contents without power ---
    ROM, or an SSD/HDD, is the example.
  \item Cache sits between the CPU and main memory and is a smaller, faster
    technology, so it is faster than RAM. It only keeps copies of data that
    already exist in RAM, so if the cache is lost the data is still in RAM;
    it is therefore not a backup.
  \item Temporal locality is the reuse of the same item soon after it was
    last accessed; spatial locality is the use of data at nearby addresses.
    Because programs show both patterns, keeping recently used and nearby
    data in the small fast cache pays off.
\end{enumerate}
\endgroup
\end{document}
